
\chapter{Neuro-Adaptive Control}

Before we go, you should distinguish between adaptive control-like NN control and Deep Learning (DL) based control.
Generally, former one employs relatively shallow NNs with few hidden layers and small number of neurons.
Also, it employs online learning scheme to update NN weights during the control process, ensuring the closed-loop stability using Lyapunov stability analysis.
The online learning stability is obtained from learning with currently measured error, which leads to local approximation of unknown dynamics.
On the other hand, DL based control generally employs deep NNs with many hidden layers and large number of neurons.
It often employs offline learning scheme to train NN weights before the control process using large dataset, which leads to semi-global approximation of unknown dynamics, in interested region.
Therefore, NAC and DL based control can be interpreted as complementary approaches rather than competing ones.
In computational point of view, former one requires less computational resources due to the small size of NNs, while latter one requires large computational resources due to the large size of NNs, thereby former one is more suitable for real-time control applications with limited computational resources and latter one is more suitable for high-performance control applications with sufficient computational resources.

\hfill

The capability of neural networks (NNs) to approximate unknown nonlinear functions are widely known in the literature, which is often referred to as the universal approximation property.
Most cited works about the capability are \cite{Barron:1993aa,Rolnick:2018aa,Kidger:2020aa}, first one for shallow NNs and the latter two for deep NNs.

Many literature have employed NNs to use their universal approximation property to approximate unknown system dynamics in adaptive control, which is often referred to as neuro-adaptive control (NAC) or neural adaptive control.
In this chapter, we review the existing NAC literature.

\section{Adaptation Laws}

Most popular adaptation law in the NAC literature is based on Lyapunov stability analysis, see \cite{Patil:2022aa}. 

\hfill

Few literature derived the adaptation law based on optimization approach.
In \cite{Esfandiari:2014aa,Esfandiari:2015aa}, Esfandiari \etal proposed an adaptation law derived from Back-Propagation (BP) algorithm to minimize the NN approximation error, while ensuring the closed-loop stability using Lyapunov direct method.
The authors used a static approximation to obtain gradient of objective function with respect to NN weights, $\pptfrac{}{\estwth}(\frac{1}{2}\mv{e}^\top\mv{e})$ while $\ddtt\mv{e}=\mv{0}$, which some authors does not agree with.

\hfill

In addition, approximation error also can be approximated using simple adaptation law, $\ddtt\widehat{\epsilon} = - \gamma_\epsilon \epsilon$, where $\gamma_\epsilon \in \R_{>0}$ is a positive constant, as in \cite{Zhao:2016aa}.
In this case, additional term $\tfrac{1}{2\gamma_\epsilon} \widetilde{\epsilon}^2$ is added to the Lyapunov function candidate to ensure the closed-loop stability including the approximation error dynamics.
In contrast, most of the literature assumes that the approximation error is bounded in a compact set (known as approximation region) without considering its dynamics.

\subsection{Boundedness of NN Weights}

In adaptive control literature, popular approaches to ensure the boundedness of NN weights are projection operator \cite{Patil:2022aa} and $\sigma$-modification \cite{He:2016aa,Gao:2006aa,Zhao:2016aa,Zhao:2025aa}, $\epsilon$-modification \cite{Esfandiari:2014aa,Esfandiari:2015aa,Zhao:2016aa}.

\section{Input Constraints}

% \begin{table}[ht]
%     \centering
%     \begin{tabular}{c|c|c|c}
%         \hline
%         \textbf{Reference} & System Type & Input Constraint Type & Auxiliary System Type \\
%         \hline
%         \hline
%         \cite{He:2016aa} & SISO & Saturation & Low-pass filter \\
%         \hline
%         \cite{Li:2024aa} & MIMO & Saturation & Smooth saturation function \\
%         \hline
%     \end{tabular}
%     \caption{Summary of Auxiliary System Based Approaches for Input Constraints in NAC Literature}
% \end{table}



\subsection{Auxiliary System Based Approach (Input Constraints)}

Auxiliary system based approach (input constraints) is widely used in the NAC literature to handle input constraints.
There are two types of auxiliary systems in the literature.

\hfill

In \cite{Esfandiari:2014aa,Esfandiari:2015aa}, Esfandiari \etal proposed a simple form of auxiliary system to handle individual actuator saturation, \ie $\ddtt\mv{\zeta} = \mm{K}\mv{\zeta}+\mv{h}\Delta\mv{u}$, $\mv{\zeta}(0) = \mv{0}$, where $\Delta \mv{u} = \mv{u} - \mysat(\mv{u})$ and $\mm{K}$ is a Hurwitz matrix.
The authors mentioned that the auxiliary system can be interpreted as $\mv{\zeta}$ represents a filtered version of $\Delta\mv{u}$ and remains zero as long as no saturation occurs, and in the presence of saturation, $\mv{\zeta}$ becomes nonzero.
More recently, similar form of the auxiliary system can be found in \cite{Yu:2022aa}.

\hfill

Besides, some literature designs relatively complex auxiliary systems to handle input saturation.
Representative form can be found in \cite{He:2016aa,Zhao:2016aa,Gao:2026aa}, where the authors designed the auxiliary system as follows:
\begin{equation}
    \ddtt \mv{\zeta} =
    \begin{cases}
        -\mm{K}\mv{\zeta}
        -
        \tfrac{
            \vert \mv{\zeta}^\top \Delta \mv{u} \vert
            + 
            \tfrac{1}{2} \Delta \mv{u}^\top \Delta \mv{u}
        }{\norm{\mv{\zeta}}^2}
        \mv{\zeta}
        +
        \Delta \mv{u}
        , & \norm{\mv{\zeta}} \ge \mu
        \\
        \mv{0}
        , & \norm{\mv{\zeta}} < \mu
    \end{cases}
    ,
\end{equation}
where $\mu > 0$ is a small positive constant.

\subsection{Smooth Saturation Function}

Some literature employs smooth saturation function to handle input saturation.
The input saturation is not directly considered in the control design, but only used in upper bounds analysis of the control error and observer error.
Since the smooth saturation function encloses the actual saturation function, control input is guaranteed to be within the allowable input space, ensuring the closed-loop stability.

\hfill

In \cite{Li:2024aa}, the input saturation of individual actuator is handled by simply employing the smooth saturation function, $\tanh(\cdot)$, in the control input design.

\subsection{Other Approaches}

As a intelligent anti-windup approach, the authors in \cite{Gao:2006aa} employed NNs to compensate saturated portion of control input, $\Delta \mv{u} := \mv{u} - \mysat(\mv{u})$.

In addition to the auxiliary system based approach, the authors in \cite{Chen:2011aa} employed a command filter to limit magnitude and rate of the virtual control input in backstepping design, ensuring the control input to be within the allowable input space.
In the auxiliary system, the difference between the virtual control input and the filtered virtual control input is considered.

In \cite{Esfandiari:2017aa}, adaptation law includes a term to handle input saturation, which is bounded even with infinite control input, \ie 
$
    \ddtt\estwth 
    = 
    -\alpha
    \left(
        (\cdot) + (\cdot)\tfrac{\norm{\Delta \mv{u}}}{\norm{\Delta \mv{u}}^2}
    \right)
$.

\subsection{Practical Applications}

In \cite{Gao:2026aa}, Gao \etal handles input saturation of each joint actuator in flexible robotic manipulator.
In \cite{Yu:2022aa}, Yu \etal handles each mobile robot's driving torque saturation in swarm control of omnidirectional mobile robots.
In \cite{Zhao:2016aa}, Zhao \etal handles input saturation of each active suspension's force actuator in vehicle active suspension control.

\section{State Constraints}

In this section, we review the existing barrier function based approaches for the NAC design with state constraints. 
There are mainly two approaches: control barrier function (CBF) approach and barrier Lyapunov function (BLF) approach.

\subsection{Preliminaries}

We first define the following system dynamics:
\begin{equation}\label{eq:sys:dyn}
    \ddtt{\mv{x}} = \mv{f}(\mv{x},\mv{u})
    ,
\end{equation}
where $\mv{x} \in \mathcal{X} \subset \mathbb{R}^n$ is the system state, $\mv{u} \in \mathcal{U} \subset \mathbb{R}^m$ is the control input, and $\mv{f}: \mathcal{X} \times \mathcal{U} \rightarrow \mathbb{R}^n$ is a locally Lipschitz continuous function representing the system dynamics.
The set $\mathcal{X}$ and $\mathcal{U}$ are the allowable state space and control input space, respectively.
Specifically, the set $\mathcal{U}$ is considered as control input constraint set.

The state is supposed to be safe within a predefined safe set $\mathcal{S} \subset \mathcal{X}$.
Consider a state feedback controller $\mv{\pi}: \mathcal{X} \rightarrow \mathcal{U}$ that renders the closed-loop system \eqref{eq:sys:dyn} safe within the safe set $\mathcal{S}$, which will be defined later.

The definitions of barrier functions and safe sets are briefly introduced in \cite{Ames:2019aa, Xiao:2019aa, agrawal:2021aa}.

\begin{definition}[Safe Set]
The safe set $\mathcal{S}$ is defined as the superlevel set of a continuously differentiable function $h: \mathcal{X} \rightarrow \mathbb{R}$ as
\begin{equation} \label{eq:safe:set}
    \begin{aligned}
        \mathcal{S} := &
        \left\{
            \mv{x} \in \mathcal{X} : h(\mv{x}) \le 0
        \right\}
        ,
        \\
        \partial \mathcal{S} := &
        \left\{
            \mv{x} \in \mathcal{X} : h(\mv{x}) = 0
        \right\}
        ,
        \\
        \myint(\mathcal{S}) := &
        \left\{
            \mv{x} \in \mathcal{X} : h(\mv{x}) < 0
        \right\}
        .
    \end{aligned}
\end{equation}
The function $h(\mv{x})$ is called a barrier function if it satisfies the following condition to ensure the safety of the system within the safe set $\mathcal{S}$.
\end{definition}

Many literature defines safe set when $h(\mv{x}) \ge 0$.
However, we define the safe set when $h(\mv{x}) \le 0$ to be consistent with the convex optimization literature.
According to reader's problem formulation, allowable state set and safe set can be subset of each other.

\begin{definition}[Barrier Function] \label{def:barrier:function}
    A continuously differentiable function $h: \mathcal{X} \rightarrow \mathbb{R}$ is said to be a barrier function for the system \eqref{eq:sys:dyn} if there exists an extended class $\mathcal{K}$ function $\alpha$ such that, for the control input $\mv{u} = \mv{\pi}(\mv{x})$, the following condition holds:
    \begin{equation}
        \ddtt h(\mv{x}) =
        \pptfrac{h(\mv{x})}{\mv{x}} \mv{f}(\mv{x}, \mv{\pi}(\mv{x})) \le -\alpha(h(\mv{x}))
        ,
        \quad
        \forall \mv{x} \in \mathcal{X}
        .
    \end{equation}
\end{definition}

By considering control input $\mv{u}$ as a variable, the barrier function condition in Definition \ref{def:barrier:function} can be rewritten as Control Barrier Function (CBF) which will be introduced in Section \ref{sec:cbf}.

\begin{theorem} [Forward Invariance of Safe Set with Barrier Function {\cite[Theorem 1]{Xiao:2019aa}}]
    Given a set $\mathcal{S}$ defined in \eqref{eq:safe:set}, if there exists a barrier function $h: \mathcal{X} \rightarrow \mathbb{R}$ for the system \eqref{eq:sys:dyn}, then the safe set $\mathcal{S}$ is forward invariant for the closed-loop system.
\end{theorem}


In \cite{Ames:2019aa}, Ames \etal 


\subsection{Control Barrier Function (CBF) Approach} \label{sec:cbf}

The control barrier function (CBF)  





\subsection{Barrier Lyapunov Function (BLF) Approach} \label{sec:blf}

Unlike the CBF approach, the barrier Lyapunov function (BLF) approach directly incorporates the safety constraints into stability analysis by constructing a Lyapunov function that diverges to infinity as the system state approaches the boundary of the safe set.

In \cite{Liu:2023aa}, 

In \cite{Zhao:2025aa}, the modified BLF approach is proposed to ensure finite resolvability of the output constraints.
Notably, the authors validated their method by imposing and removing the constraints repeatedly.

\section{Remaining Topics}

\subsection{Extension to Output Feedback Case}

High-gain observer \cite{Khalil:2002aa} is employed in \cite{Esfandiari:2015aa,He:2016aa}, which offers separation principle between the controller and observer designs, \ie controller's performance recovers to the full-state feedback case as the observer gain increases.
The authors in \cite{Esfandiari:2015aa} mentioned that due to the high-gain nature of the observer, peaking phenomenon should be carefully handled to ensure that the system states remain within the allowable state space.

\hfill

Notably, NNs can be employed in the observer design to estimate unmeasured states, as well.
In \cite{Zhao:2016aa}, the NN-based observers can be found, but the authors did not provide closed-loop stability analysis including the observer dynamics.

\hfill 

Additionally, in \cite{Esfandiari:2017aa}, state identification approach can be found.

\subsection{Parameter Selection Guidelines}

Due to the trade off between control performance and stability, parameter selection guidelines are important in practical applications.
Most of the literature do not provide clear parameter selection guidelines of NNs.

\hfill

Particle swarm optimization (PSO) based approach can be found in \cite{Zhao:2016aa}, where the authors optimized NN parameters (not including structure) using PSO algorithm.

\subsection{For General Systems}

Most of the existing NAC literature focus on specific class of systems, such as strict-feedback systems.
There are few literature that consider more general class of systems.

\hfill

Regarding MIMO systems, most of literature assume that control effectiveness matrix is known \cite{Yu:2022aa}, its upper and lower bounds are known \cite{}, or it is nonsingular \cite{}.
In \cite{Chen:2011aa}, Chen \etal relaxed widely used assumption that control effectiveness matrix is nonsingular, by considering spectral radius of the matrix, see Lemma \ref{appx:lem:spectral:radius} in Appendix.

More generally, in \cite{Esfandiari:2017aa}, Esfandiari \etal proposed a NAC design for nonaffine nonlinear systems, with assumptions that it is stabilizable, $f(0,0)=0$ and Lipschitz continuous in a compact set containing the origin.

\subsection{Reinforcement Learning Approach}

Considering the objective function in the form of integral of instantaneous cost, \ie $J = \int_0^\infty c(\mv{x},\mv{x}_d)\,\der t$, reinforcement learning (RL) approach based on optimal control theory can be found in \cite{Esfandiari:2017aa}.